% from ega3-7-fr.tex
\oldpage[III]{190}
\subsection{Cas des anneaux locaux noethériens}
\label{subsection:III.7.5.fr}

\begin{env}[7.5.1]
Soient $A$ un anneau local noethérien, $\mathfrak m$ son idéal maximal, et pour
tout $A$-module $M$, désignons par $\widehat M$ son séparé complété pour la
topologie $\mathfrak m$-préadique, isomorphe à
\[
  \varprojlim_n\bigl(M\otimes_A(A/\mathfrak m^{n+1})\bigr)
  =
  \varprojlim_n\bigl(M/\mathfrak m^{n+1}M\bigr).
\]
Soit $T$ un foncteur covariant additif de $\mathbf{Ab}_A$ dans $\mathbf{Ab}$;
les homomorphismes canoniques (7.2.3.1)
\[
  T(M)\otimes_A(A/\mathfrak m^{n+1})
  \longrightarrow
  T\bigl(M\otimes_A(A/\mathfrak m^{n+1})\bigr)
\]
forment évidemment un système projectif de $A$-homomorphismes, qui donnent donc
à la limite un $\widehat A$-homomorphisme fonctoriel en $M$
\[
  \widehat{T(M)}
  \longrightarrow
  \varprojlim_n T(M_n)
  \tag{7.5.1.1}\label{III.7.5.1.1.fr}
\]
où l'on a posé
\[
  M_n=M\otimes_A(A/\mathfrak m^{n+1}),
  \qquad
  A_n=A/\mathfrak m^{n+1}.
\]
\end{env}

\begin{proposition}[7.5.2]
Soient $A$ un anneau local noethérien d'idéal maximal $\mathfrak m$,
$k=A/\mathfrak m$ son corps résiduel, $T$ un foncteur covariant additif de
$\mathbf{Ab}_A$ dans $\mathbf{Ab}$, semi-exact et commutant aux limites
inductives. On suppose en outre que pour tout $A$-module de type fini $M$,
$T(M)$ est un $A$-module de type fini et que l'homomorphisme canonique
(7.5.1.1) est un isomorphisme. Sous ces conditions, les propriétés suivantes
sont équivalentes :
\begin{enumerate}
  \item[a)] $T$ est exact à droite.
  \item[b)] Pour tout $n$, le foncteur $N\mapsto T(N)$ est exact à droite dans
    la catégorie des $A_n$-modules de type fini (ce qui revient à dire que $T$
    est exact à droite dans la catégorie des $A$-modules de longueur finie).
  \item[c)] L'homomorphisme canonique
    \[
      T(\varepsilon):T(A)\longrightarrow T(k)
    \]
    est surjectif.
  \item[d)] Pour tout $n$ assez grand, l'homomorphisme canonique
    \[
      T(A_n)\longrightarrow T(k)
    \]
    est surjectif.
\end{enumerate}
\end{proposition}

\begin{proof}
Il est clair que a) entraîne b). Montrons que b) entraîne a), c'est-à-dire que
si $u:M\to N$ est un épimorphisme de $A$-modules, $T(u)$ est un épimorphisme.
Comme $T$ commute aux limites inductives et que le foncteur $\varinjlim$ est
exact dans la catégorie des modules (pour des ensembles d'indices filtrants),
on peut se borner au cas où $M$ et $N$ sont de type fini. Comme $T(M)$ et
$T(N)$ sont alors de type fini, et que $A$ est un anneau local noethérien, il
suffit de montrer que
\[
  \widehat{T(u)}:
  \widehat{T(M)}
  \longrightarrow
  \widehat{T(N)}
\]
est surjectif ($0_{\mathrm I}$, 7.3.5 et $0_{\mathrm I}$, 6.4.1). Par hypothèse,
$\widehat{T(M)}$ et $\widehat{T(N)}$ sont respectivement $\varprojlim_n T(M_n)$
et $\varprojlim_n T(N_n)$, donc $\widehat{T(u)}$ est limite du système projectif
d'homomorphismes
\[
  T(u\otimes1_{A_n}):T(M_n)\longrightarrow T(N_n).
\]
Or, b) signifie que ces homomorphismes sont surjectifs; en outre, $T(M_n)$ est
un $A_n$-module de type fini, et $A_n$ est un anneau artinien par hypothèse; on
en conclut que $\widehat{T(u)}$ est surjectif ($0_{\mathrm I}$, 13.1.2 et
13.2.2). Il est clair que a) entraîne c), et comme $T(\varepsilon)$ se factorise
en
\[
  T(A)\longrightarrow T(A_n)\longrightarrow T(k),
\]
c) implique d); enfin, il résulte de (7.2.7) que b) et d) sont équivalentes
puisque $T$ est semi-exact dans $\mathbf{Ab}_{A_n}$, ce qui achève la
démonstration.
\end{proof}

\begin{corollary}[7.5.3]
Sous les hypothèses générales de (7.5.2), si $T(k)=0$, alors $T(M)=0$ pour tout
$A$-module $M$.
\end{corollary}

\begin{proof}
